\documentclass[main.tex]{subfiles}
\usepackage{algorithm}
\usepackage{algpseudocode}
\begin{document}

%%%%%%%%%%%%%%%%%%%%%%%%%%%%%%%%%%%%%%%%%%%%%%%%%%%%%%%%%%%%%%%%%%%%%%%%%%%%%%%%%%%%%%%
%%%%%%%%%%%%%%%%%%%%%%%%%%%%%%%%%%%%%%%%%%%%%%%%%%%%%%%%%%%%%%%%%%%%%%%%%%%%%%%%%%%%%%%
%%%%%%%%%%%%%%%%%%%%%%%%%%%%%%%%%%%%%%%%%%%%%%%%%%%%%%%%%%%%%%%%%%%%%%%%%%%%%%%%%%%%%%%
%%%%%%%%%%%%%%%%%%%%%%%%%%%%%%%%%%%%%%%%%%%%%%%%%%%%%%%%%%%%%%%%%%%%%%%%%%%%%%%%%%%%%%%
%%%%%%%%%%%%%%%%%%%%%%%%%%%%%%%%%%%%%%%%%%%%%%%%%%%%%%%%%%%%%%%%%%%%%%%%%%%%%%%%%%%%%%%
\subsection*{Notations}
\label{notations}

In this section, we introduce the objects we are going to work with and the summary of the notation in \ref{tab:standard_definitions} that are going to hold throughout the work.

%%%%%%%%%%%%%%%%%%%%%%%%%%%%% People
We call Users the set of people under study we index them with greek letters ($\alpha$,$\beta$).

%%%%%%%%%%%%%%%%%%%%%%%%%%%%% Geometry
We call grid the geometry of squared $N_{grid}$ cells  defined by geohash 7 tiling with side size $L_{cell} = 70$ m, we index each cell with roman letters i,j,k. 
We therefore denote a single cell of the grid as $g_i$.

%%%%%%%%%%%%%%%%%%%%%%%%%%%%% Stop point
A stop point ($sp_t^{\alpha}$) is a couple of coordinates  $(lat^{\alpha}_t,lon^{\alpha}_t)$ and a starting (t) and ending time ($\hat{t}$) such that the mobile phone's signal remains within a cell of the grid for a minimum duration of $t^{sp}$ 10 minutes: $\hat{t} - t > t^{sp}$. We consider the time of the stop point as the beginning of the stop.
We denote stay time $\delta t_{t} = \hat{t} - t$ and is variable in a trajectory.

%%%%%%%%%%%%%%%%%%%%%%%%%%%%% Trajectories
To each user it is associated a trajectory. A trajectory  $traj^\alpha$ is a succession of $N^{\alpha}$ stop points distributed in time $T^{\alpha} = \{t_{\tau}: \tau = 0,1,..,N^{\alpha}\}$ and the time duration of the trajectory is $|T^{\alpha}| = t_{N^{\alpha}}-t_0$. We define trajectories in three equivalent ways as,
%
\begin{itemize}
    \item $traj^{\alpha} = \{sp^{\alpha}_t\}_{t \in T^{\alpha}}$ (succession of lat lon)
    \item $traj^{\alpha} = \{g^{\alpha}_{it}\}_{t \in T^{\alpha}}$ (succession of cells)
    \item $traj^{\alpha} = \{(x^{\alpha}_t,y^{\alpha}_t)\}_{t \in T^{\alpha}}$ (succession of points in the tangent space).
\end{itemize}
%
We call $Trajs = \{traj_{\alpha}\}_{\alpha \in Users}$ the set of trajectories.

%%%%%%%%%%%%%%%%%%%%%%%%%%%%% time intervals and split into epochs
Each trajectory is characterized with an interval $[t_{start}^{\alpha},t_{end}^{\alpha}]$ that is contained in the overall period of study $[T_{start} = 15/01/2020,T_{end} = 01/09/2020]$.

%%%%%%%%%%%%%%%%%%%%%%%%%%%%% covid epochs
We futher divide the overall period into three different periods representing different stages of lockdown restrictions that we index with suffixes $\{pre,during,post\}$ and generically referred to with suffix p.

This partition is defined by the time intervals $\{[T^{pre}_{start} = 15/01/2020,T^{pre}_{end} = 15/03/2020],[T^{during}_{start} = 15/03/2020,T^{during}_{end} = 15/05/2020],[T^{post}_{start} = 15/05/2020,T^{post}_{end} = 01/09/2020]\}$ and induces a partition in the set of users that we call $\{Users^{pre},Users^{during},Users^{post}\}$ with the condition for a user to belong to a certain class is that $t_{end}^{\alpha} \le T^{p}_{end}$.

In the case user $\alpha$ spans multiple periods, the trajectory is cut into multiple sub-trajectories, and each sub-trajectory is considered for the different partition $Users^{p}$.

%%%%%%%%%%%%%%%%%%%%%%%%%%%%% Filtering
We also define a time burst filter $t^{burst}$ as the minimum time interval within two successive stop points.

%%%%%%%%%%%%%%%%%%%%%%%%%%%%% Distance
We define as the geodesic distance $d_g(\Vec{x},\Vec{y})$ between two points $\Vec{x}$,$\Vec{y}$.

%%%%%%%%%%%%%%%%%%%%%%%%%%%%% Mobility network
In the end we define the mobility network of user $\alpha$ as $G^{\alpha} = (N^{\alpha},E^{\alpha})$.
This is a weighted directed graph where the nodes (n) are the stop points, while the edges (e) are the transitions between successive stop points and the weight $w_{e}$ is the number of time such transition happened . The graph is naturally loop-less. We call $k^{\alpha,in}_n$ the number of time node n is visited.
%
\begin{table}[h!]
    \centering
    \begin{tabular}{|c|c|}
        \hline
        % grid
        $grid$, (i,j,k) &  set of cells geohash 7, (indices) \\
        \hline
        $d_{g}(\Vec{x},\Vec{y})$ &  geodesic distance between two points $\Vec{x}$ and $\Vec{y}$\\
        \hline
        % Users & Trajs (not conditioned)
        Users & Set of people \\
        \hline
        $Trajs$ &  Set of all trajectories \\
        \hline
        $(\alpha,\beta)$ & index user \\
        \hline
        % Trajectories components
        $l^{\alpha}$ & length of the trip user $\alpha$ \\
        \hline
        % Time 
        $T^{\alpha}$ & set of times for each stop point for trajectory $\alpha$\\
        \hline
        $|T^{\alpha}|$ & time duration for trajectory $\alpha$\\
        \hline

        $T_{start}$,$T_{end}$ & start and end time analysis \\
        \hline
        $T^{p}_{start}$,$T^{p}_{end}$ & start and end time of period p \\
        \hline
        $t_{start}^{\alpha}$,$t_{end}^{\alpha}$ & start and end time user $\alpha$ \\
        \hline
        $\tau$ & integer index time of a stop point in a trajectory \\
        \hline
        % filters time
        $t^{sp} (=10 min)$ & interval definition stop point\\
        \hline
        $\delta t_{t}$ & duration of $sp_{t}$\\
        \hline
        % Users and Trajs (condtioned)
        $Users^{p}$ & Set of people per period p \\
        \hline
        $Trajs^{p}$ &  Set of all trajectories per period p \\
        \hline
    \end{tabular}
    \caption{Standard Definitions}
    \label{tab:standard_definitions}
\end{table}
%

%%%%%%%%%%%%%%%%%%%%%%%%%%%%%%%%%%%%%%%%%%%%%%%%%%%%%%%%%%%%%%%%%%%%%%%%%%%%%%%%%%%%%%%
%%%%%%%%%%%%%%%%%%%%%%%%%%%%%%%%%%%%%%%%%%%%%%%%%%%%%%%%%%%%%%%%%%%%%%%%%%%%%%%%%%%%%%%
%%%%%%%%%%%%%%%%%%%%%%%%%%%%%%%%%%%%%%%%%%%%%%%%%%%%%%%%%%%%%%%%%%%%%%%%%%%%%%%%%%%%%%%
%%%%%%%%%%%%%%%%%%%%%%%%%%%%%%%%%%%%%%%%%%%%%%%%%%%%%%%%%%%%%%%%%%%%%%%%%%%%%%%%%%%%%%%
%%%%%%%%%%%%%%%%%%%%%%%%%%%%%%%%%%%%%%%%%%%%%%%%%%%%%%%%%%%%%%%%%%%%%%%%%%%%%%%%%%%%%%%
\subsection*{Data}
\label{dataset}

Cuebiq Inc. provided the human mobility data used here for research purposes. This data was collected from anonymous mobile phone users who have opted-in to give access to their location data anonymously, through a GDPR-compliant framework. Researchers queried the mobility data through an auditable, cloud-hosted sandbox environment, receiving aggregate outputs in return. The datasets contain records of people living in California and Massachusetts from January 2020 to the end of September 2020. The former dataset is described in Tab.\ref{tabPeopleCal}. Given the unique nature of the data used in this study (enhanced mobile phone data collected by Cuebiq, characterized by concentrated bursts of data collection within short time windows), it was necessary to implement a burst filtering process in order to ensure consistency in the acquisition of positional data. Specifically, the filter imposes that $\forall traj_{\alpha} \in Trajs, \forall sp_{t_{\tau}},sp_{t_{\tau + 1}} \in traj_{\alpha} $ that $t_{\tau + 1} - t_{\tau} > t^{burst}$, that is, two successive stop points are further apart in time by at least $t^{burst}$. In this discussion, $t^{burst} = 1$ hour; however, alternative threshold values of 8 and 24 hours were also applied and are available in the Appendix for reference. We note here that the duration time $\delta t_{t_{\tau}}$ for a stop point is not the time between two successive stop points $sp_{t_{\tau}},sp_{t_{\tau + 1}}$ after the introduction of the burst filter.

%%%%%%%%%%%%%%%%%%%%%%%%%%%%%%%%%%%%%%%%%%%%%%%%%%%%%%%%%%%%%%%%%%%%%%%%%%%%%%%%%%%%%%%
%%%%%%%%%%%%%%%%%%%%%%%%%%%%%%%%%%%%%%%%%%%%%%%%%%%%%%%%%%%%%%%%%%%%%%%%%%%%%%%%%%%%%%%
%%%%%%%%%%%%%%%%%%%%%%%%%%%%%%%%%%%%%%%%%%%%%%%%%%%%%%%%%%%%%%%%%%%%%%%%%%%%%%%%%%%%%%%
%%%%%%%%%%%%%%%%%%%%%%%%%%%%%%%%%%%%%%%%%%%%%%%%%%%%%%%%%%%%%%%%%%%%%%%%%%%%%%%%%%%%%%%
%%%%%%%%%%%%%%%%%%%%%%%%%%%%%%%%%%%%%%%%%%%%%%%%%%%%%%%%%%%%%%%%%%%%%%%%%%%%%%%%%%%%%%%
\subsection*{Metrics}

We examine various measures that emphasize distinct types of comprehension, including geometric, information-theoretic, and visitation patterns. For starters, we have used a modified version of the radius of gyration, shown in Eq.~\ref{rg}. This measures the spatial spread of a person's trajectory around its time-weighted centre of mass: the longer a user dwells at a location, the greater its contribution to the centroid and hence to the radius. The time-weighted centre of mass $\Vec{x^{\alpha}_{cm}}$ is defined as
%
\begin{equation}
    \Vec{x^{\alpha}_{cm}} = \frac{\sum_{t \in T^{\alpha}} \delta t_{t}\, \Vec{x^{\alpha}_{t}}}{\sum_{t \in T^{\alpha}} \delta t_{t}},
    \label{cm}
\end{equation}
%
and the radius of gyration is then
%
\begin{equation}
     R^{\alpha}_{g} = \sqrt{\frac{\displaystyle\sum_{t \in T^{\alpha}} \delta t_{t}\; d_{g}(\Vec{x^{\alpha}_{t}},\Vec{x^{\alpha}_{cm}})^2}{\displaystyle\sum_{t \in T^{\alpha}} \delta t_{t}}}
     \label{rg}
\end{equation}
%
where $\delta t_{t}$ is the time spent at the stop point of coordinates $\Vec{x^{\alpha}_{t}}$, $\Vec{x^{\alpha}_{cm}}$ is the time-weighted centre of mass of the full trajectory (Eq.~\ref{cm}), and $d_{g}$ is the geodesic distance projected onto the local tangent plane. All distances are expressed in kilometres. The centre of mass is computed once over the entire observation period for each user--period pair; it is not updated incrementally as new stops are observed. Another quantity that expresses geometric information is the distribution of stop points in trajectories in their intrinsic reference frame, re-scaled by the variances on the principal axes \cite{radius_gyration}. For each user $\alpha$, the procedure is:
%
\begin{enumerate}
    \item Project all stop coordinates to a local tangent plane centred at the trajectory's mean position, yielding $(x_t^{\alpha}, y_t^{\alpha})$ in metres.
    \item Centre the point cloud: $\tilde{x}_t = x_t^{\alpha} - \bar{x}^{\alpha}$, $\tilde{y}_t = y_t^{\alpha} - \bar{y}^{\alpha}$.
    \item Compute the second-moment (inertia) tensor:
    \[
        I_{xx} = \sum_t \tilde{x}_t^2,\quad
        I_{yy} = \sum_t \tilde{y}_t^2,\quad
        I_{xy} = \sum_t \tilde{x}_t\tilde{y}_t.
    \]
    \item Find the principal rotation angle $\theta$ that diagonalises the tensor:
    \[
        \mu = \sqrt{4I_{xy}^2 + (I_{xx}-I_{yy})^2}, \qquad
        \cos\theta = \frac{-I_{xy}}{\left(\frac{I_{xx}-I_{yy}+\mu}{2}\right)
                      \sqrt{1 + \frac{I_{xy}^2}{\left(\frac{I_{xx}-I_{yy}+\mu}{2}\right)^2}}}.
    \]
    \item Rotate to the principal frame: $(x'_t, y'_t) = R_\theta\,(\tilde{x}_t, \tilde{y}_t)$.
    \item Compute the root-mean-square spreads: $\sigma_{x'} = \sqrt{\langle x'^2\rangle}$, $\sigma_{y'} = \sqrt{\langle y'^2\rangle}$.
    \item Normalise: $\hat{x}_t = x'_t/\sigma_{x'}$, $\hat{y}_t = y'_t/\sigma_{y'}$.
\end{enumerate}
%
The resulting 2D probability distribution on $(\hat{x},\hat{y})$ is the Gonzalez trajectory shape, which is independent of the trajectory's overall size and orientation.
\label{phi}
%
We also measured the curve of newly visited places in time and the average frequency of the ten most frequented places per person, shown in figure \ref{visitation increase-frequency}.
Let $r$ denote the rank $r = 1,\ldots,10$. The normalised visit frequency at rank $r$ in period $p$ is given by Eq.~\ref{k_avg}.
%
\begin{equation}
    \langle k \rangle^{r,p} =  \left\langle \frac{k^{\alpha,r}}{\sum_{r' = 1}^{\min(10,\,r_{\max})}k^{\alpha,r'}} \right\rangle_{{\alpha} \in {Trajs^p}}
    \label{k_avg}
\end{equation}
%
Here, \( p \) denotes the time period, \( r \) represents the rank of a node in an individual’s mobility network (with \( r = 1 \) being the most visited location), and \( k^{\alpha, r} \) is the number of trips made by individual \( \alpha \) to the node of rank \( r \). The distribution is computed separately for each period \( p \), considering only the individuals active during that period.

This metric captures how the frequency of visits to the most important locations changes across the population over different time periods.
In the end, we calculated entropic measures that quantify the degree of uncertainty of the trajectories in three different ways (Eqs.~\ref{rdm_entropy}--\ref{real_entropy}).
%
\begin{equation}
    H_{\text{random}} = \log_2 n \label{rdm_entropy}
\end{equation}
%
\begin{equation}
    H_{\text{uncorrelated}} = -\sum_{i=1}^n p_i \log_2 p_i
    \label{unc_entropy}
\end{equation}
%
\begin{equation}
    H_{\text{real}} = \left(\frac{1}{T}\sum_{t=1}^{T} \Lambda_t\right)^{-1} \log_2 T
    \label{real_entropy}
\end{equation}
%
The random entropy (Eq.~\ref{rdm_entropy}) depends only on the number $n$ of distinct visited locations. The uncorrelated entropy (Eq.~\ref{unc_entropy}) adds sensitivity to the relative frequency of visits $p_i$, treating successive stops as independent. The real entropy (Eq.~\ref{real_entropy}) is a Lempel-Ziv-78 estimator that accounts for the sequential order of visits and converges to the true entropy rate of the trajectory \cite{Lempel-Ziv, entropic_measures}. Here $T$ is the total number of stop points and $\Lambda_t$ is the length of the shortest substring starting at position $t$ that has not previously appeared in positions $1, \ldots, t-1$. These three measures provide a progressively finer view of trajectory complexity. The $k$-radius of gyration extends Eq.~\ref{rg} by restricting the computation to the $k$ most-visited locations (by visit count), capturing how much of a user's mobility is concentrated near their habitual places.



\end{document}

